\documentclass[12pt]{article}
\newcommand{\bottomMargin}{2cm}


\newcommand{\header}{
	ДИЗАЙН И АНАЛИЗ НА АЛГОРИТМИ - \\ ПРАКТИКУМ\\
	\vspace{0.1cm}
        Летен семестър, 2026 г., първо контролно\\
	\vspace{0.1cm}
}
\newcommand{\tl}{$0.1$ сек.}
\newcommand{\ml}{$256$ MB}

\input{structure.tex}

\begin{document}

\section{Задача К2. МинМакс}

Сашка разглежда любимата си редица $a_1, a_2, \dots, a_N$ от $N$ цели положителни
числа, но нещо не и харесва. Тя забелязва, че разликата между най-голямото и най-малкото число в редицата е твърде голяма, заради това тя иска да я промени. Тя е съгласна до $S$ пъти да изтрие някое число $a_x$ от редицата и да напише $a_x - 1$ или
$a_x + 1$ на негово място. Например при $a = (1, 3, 2)$, Сашка може с три триения да я
превърне в $a = (2, 1, 2)$.

Напишете програма \textbf{minmax}, която да намира минималната постижима разлика между най-големия и най-малкия елемент от редицата след най-много $S$ триения.

\subsection{Вход}
От първия ред на стандартния вход се въвеждат $2$ цели естествени числа $N$ и $S$ -- съответно броят на числата в редицата и позволеният брой триения. На втория ред се въвеждат $N$ числа -- $a_1, a_2, \dots, a_N$ -- числата от редицата на Сашка.

\subsection{Изход}
На единствения ред на стандартния изход изведете едно цяло число -- минималната постижима разлика между
най-големия и най-малкия елемент от редицата.


\subsection{Ограничения}
\vspace{0.1em}
\begin{itemize}
	\item $1 \leq N \leq 10^5$
	\item $1 \leq S\leq 10^{14}$
    \item $1 \leq a_i \leq 10 ^9$
\end{itemize}

\subsection{Оценяване}
\noindent В $10\%$ от тестовете, всички елементи от редицата са различни и $a_i \leq N$. \\
В други тестове, носещи около $20\%$ от точките: $a_i \leq 10^6$.\\
В други тестове, носещи около $20\%$ от точките: $a_i \leq 5000, N \leq 1000$.\\
В други тестове, носещи около $30\%$ от точките: $N \leq 2000$.\\
В останалите тестове, носещи около $20\%$ от точките, няма допълнителни
ограничения.

\subsection{Пример}
\begin{table}[H]
	\begin{tblr}{|l|l|X[j]|}
		\hline
		\textbf{Вход} & \textbf{Изход} & \textbf{Обяснение на примера} \\
		\hline
		\texttt{\makecell[lt]{8 5 \\
            4 2 3 2 4 3 4 3 
		}}
		& 
		\texttt{\makecell[lt]{0}}
		& 
		{Всички елементи стават равни на $3$.}\\
		\hline
		\texttt{\makecell[lt]{7 3 \\
1 9 2 9 10 2 10 
		}}
		& 
		\texttt{\makecell[lt]{7}}
		& 
		{След оптимално прилагане на операциите,
редицата може да изглежда така:\\
$2 \ 9 \ 2 \ 9 \ 9 \  2 \ 9$}\\
		\hline
		\texttt{\makecell[lt]{7 15 \\
1 9 2 9 10 2 10
		}}
		& 
		\texttt{\makecell[lt]{3}}
		& 
		{След оптимално прилагане на операциите,
редицата може да изглежда така:\\
$6 \ 9 \ 6 \ 9 \ 9 \  6 \ 9$}\\
		\hline
	\end{tblr}
\end{table}
\FloatBarrier

\end{document}
